\documentclass[11pt,letterpaper]{article}
\usepackage[margin=1in]{geometry}
\usepackage{graphicx}
\usepackage{hyperref}
\usepackage{booktabs}
\usepackage{listings}
\usepackage{xcolor}
\usepackage{amsmath}

\definecolor{codebg}{HTML}{F5F5F5}
\lstset{
    backgroundcolor=\color{codebg},
    basicstyle=\ttfamily\small,
    breaklines=true,
    frame=single,
    rulecolor=\color{gray!40},
    numbers=left,
    numberstyle=\tiny\color{gray},
    tabsize=4,
}

\title{\textbf{Adaptive RDMA Optimization for Distributed LLM Workloads}\\[0.5em]
\large Update 9: WFA Tensor Classification Validation on LLM4Decompile Inference}
\author{Nathan Gopee\\
MS Computer Science, SUNY New Paltz\\
\texttt{gopeen1@newpaltz.edu}\\[1em]
\small Advisors: Dr.\ Wafi Danesh \& Dr.\ Ashley Suchy}
\date{February 2026}

\begin{document}
\maketitle

\begin{abstract}
\noindent
\textbf{Project Purpose.}
This update validates the Weighted Finite Automaton (WFA) tensor classifier from the adaptive RDMA middleware (Update~10) against a real workload: LLM4Decompile, an open-source LLM that decompiles x86-64 binaries back into C source code. By instrumenting KV cache access patterns during decompilation inference on a Google Colab H100 GPU in BF16 with FlashAttention~2, we capture per-layer access frequencies and run the WFA classifier to categorize each layer's KV cache as hot, warm, or cold.

Additionally, we implement SAE (Sparse Autoencoder) feature steering---the same technique behind Golden Gate Claude and the Eiffel Tower Llama demo---to demonstrate gradient-free behavior modification and measure its impact on WFA classification. Steering specific transformer layers changes their access patterns, which the WFA detects and would route to dedicated RDMA Queue Pairs.

No RDMA hardware is required---the classifier and steering engine are pure Python. The experiments demonstrate that (1)~the WFA correctly differentiates access patterns across transformer layers during autoregressive generation, and (2)~SAE steering creates observable differences in per-layer compute load that inform RDMA routing decisions.
\end{abstract}

\tableofcontents
\newpage

%=============================================================================
\section{LLM4Decompile Overview}
\label{sec:llm4decompile}
%=============================================================================

LLM4Decompile~\cite{llm4decompile} is the first open-source large language model dedicated to binary decompilation. Given x86-64 assembly (compiled with GCC at optimization levels O0--O3), the model generates human-readable C source code. Two model families exist:

\begin{itemize}
    \item \textbf{LLM4Decompile-End}: Directly decompiles assembly to C.
    \item \textbf{LLM4Decompile-Ref}: Refines Ghidra pseudo-code output into cleaner C.
\end{itemize}

\begin{table}[htbp]
\centering
\caption{LLM4Decompile model variants and re-executability scores}
\label{tab:llm4decompile-models}
\begin{tabular}{@{}lccl@{}}
\toprule
\textbf{Model} & \textbf{Parameters} & \textbf{Re-executability} & \textbf{Approach} \\
\midrule
llm4decompile-1.3b-v1.5 & 1.3B & 27.3\% & Direct (End) \\
llm4decompile-6.7b-v1.5 & 6.7B & 45.4\% & Direct (End) \\
llm4decompile-6.7b-v2 & 6.7B & 52.7\% & Ghidra refinement (Ref) \\
llm4decompile-9b-v2 & 9B & 64.9\% & Ghidra refinement (Ref) \\
llm4decompile-22b-v2 & 22B & 63.6\% & Ghidra refinement (Ref) \\
\bottomrule
\end{tabular}
\end{table}

Re-executability measures whether the decompiled C code compiles and passes the original test assertions---a strict correctness metric. We use the 6.7B v1.5 model on a Colab H100 VM in BF16 with FlashAttention~2 for maximum throughput.

\subsection{Relevance to Thesis}

LLM4Decompile connects to the thesis in two ways:

\begin{enumerate}
    \item \textbf{WFA validation workload}: Decompilation inference exercises the full autoregressive generation pipeline (prefill + decode), producing real KV cache access patterns for WFA classifier validation.
    \item \textbf{RISC-V assembly generation} (Section~14.2 of Update~10): The same model architecture could be fine-tuned on RISC-V binaries using the heterogeneous training pipeline (V100 FP16 forward/backward, Neumann FP32 master weights, 5070~Ti BF16 validation).
\end{enumerate}

%=============================================================================
\section{Experiment: Instrumenting KV Cache Access Patterns}
\label{sec:experiment}
%=============================================================================

\subsection{Setup}

\begin{table}[htbp]
\centering
\caption{Experiment environment}
\label{tab:experiment-setup}
\begin{tabular}{@{}ll@{}}
\toprule
\textbf{Component} & \textbf{Specification} \\
\midrule
Platform & Google Colab (H100 VM) \\
GPU & NVIDIA H100 80 GB (compute 9.0, Hopper) \\
Model & LLM4Binary/llm4decompile-6.7b-v1.5 \\
Precision & BF16 (native H100 Tensor Core format) \\
Attention & FlashAttention 2 \\
Framework & HuggingFace Transformers + PyTorch \\
\bottomrule
\end{tabular}
\end{table}

\subsection{Method}

We instrument the model by hooking each transformer layer's self-attention forward pass. Every time a layer's KV cache is accessed during generation, we record:

\begin{itemize}
    \item The layer index
    \item A nanosecond-resolution timestamp
    \item The cumulative access count
\end{itemize}

The WFA classifier then processes these traces, classifying each layer's KV cache stream as hot, warm, or cold based on access frequency and idle time---exactly the same logic that would run in the RDMA middleware on the two-tower cluster.

The Colab notebook (\texttt{Update9/llm4decompile\_wfa.ipynb}) implements the full pipeline. The key components are extracted below for reference; all runnable code is in the notebook.

\subsection{Instrumentation Approach}

Each transformer layer's \texttt{self\_attn.forward} method is wrapped with a hook that increments an access counter and appends a timestamp before calling the original function. This captures the real access pattern without modifying the model's computation.

During autoregressive generation, the prefill phase processes all input tokens in one forward pass (each layer accessed once), followed by the decode phase where each new token requires one forward pass through all layers. The WFA classifier observes this as a burst of rapid sequential accesses across all layers during decode---exactly the pattern that should be classified as HOT.

\subsection{WFA Classification}

The classifier uses three parameters:
\begin{itemize}
    \item $\theta_1 = 5$: warm threshold (accesses per observation window)
    \item $\theta_2 = 20$: hot threshold
    \item $t_{\text{idle}} = 0.5$s: idle timeout for demotion
\end{itemize}

For a 512-token generation on a 32-layer model, each layer's KV cache is accessed $\sim$512 times (once per generated token). All layers should classify as HOT, which validates that the classifier correctly identifies the decode-phase access pattern.

More interesting classification emerges with:
\begin{itemize}
    \item \textbf{Batch inference}: Multiple sequences with different lengths create variable access frequencies per layer.
    \item \textbf{Speculative decoding}: Draft model layers are accessed more frequently than verification layers.
    \item \textbf{Early exit}: Some layers may be skipped for easy tokens, creating warm/cold patterns.
\end{itemize}

%=============================================================================
\section{Expected Results}
\label{sec:results}
%=============================================================================

\subsection{Access Pattern Profile}

For single-sequence autoregressive generation of $N$ tokens on a model with $L$ layers:

\begin{equation}
    \text{Total KV accesses} = L \times (1 + N) = L \times (\text{prefill} + \text{decode steps})
\end{equation}

With $L = 32$ and $N = 512$: $32 \times 513 = 16{,}416$ total accesses, uniformly distributed across layers. The access frequency per layer ($\sim$513) far exceeds $\theta_2 = 20$, so all layers classify as HOT---correct for single-sequence decode.

\subsection{Phase Detection}

The \texttt{detect\_phase} function distinguishes:
\begin{itemize}
    \item \textbf{Prefill}: tokens\_per\_seq $> 8$ (assembly input is typically 50--200 tokens)
    \item \textbf{Decode}: tokens\_per\_seq $= 1$ (autoregressive generation)
\end{itemize}

The phase transition from prefill to decode is the critical moment where the RDMA middleware would switch from bulk-transfer mode (cold path, large input tensor) to streaming mode (hot path, small per-token KV updates).

\subsection{Validation Criteria}

\begin{table}[htbp]
\centering
\caption{WFA classification validation targets}
\label{tab:validation-targets}
\begin{tabular}{@{}lll@{}}
\toprule
\textbf{Scenario} & \textbf{Expected Classification} & \textbf{Target Accuracy} \\
\midrule
Single-seq decode (all layers) & HOT & 100\% \\
Prefill phase detection & Correctly identified & 100\% \\
Idle layers (simulated skip) & COLD after $t_{\text{idle}}$ & $>$95\% \\
Variable-frequency batch & Mixed HOT/WARM/COLD & $>$90\% \\
\bottomrule
\end{tabular}
\end{table}

%=============================================================================
\section{Connection to RDMA Middleware}
\label{sec:rdma-connection}
%=============================================================================

The classification data from this experiment directly informs RDMA routing decisions on the two-tower cluster:

\begin{table}[htbp]
\centering
\caption{Mapping WFA classifications to RDMA routing actions}
\label{tab:wfa-rdma-mapping}
\begin{tabular}{@{}llll@{}}
\toprule
\textbf{Classification} & \textbf{Tensor Type} & \textbf{RDMA Path} & \textbf{QP Config} \\
\midrule
HOT & KV cache (decode) & Dedicated low-latency QPs & Polled completions \\
WARM & Activations (pipeline) & Shared QPs & Event-driven \\
COLD & Weights, gradients & Batched cold channel & Overlapped with compute \\
\bottomrule
\end{tabular}
\end{table}

On the Colab T4, the classification runs locally (no network transfer). On the two-tower cluster, the same classifier output drives QP selection: a tensor classified as HOT during decode goes through a dedicated RDMA Write with Immediate on a priority Queue Pair, achieving sub-5$\mu$s transfer. A COLD tensor (gradient accumulation) gets batched and sent when 8 tensors accumulate, overlapping with the next layer's compute.

%=============================================================================
\section{SAE Feature Steering Experiment}
\label{sec:sae-steering}
%=============================================================================

Sparse Autoencoder (SAE) feature steering is a gradient-free technique for modifying LLM behavior at inference time. Originally demonstrated by Anthropic's Golden Gate Claude~\cite{golden-gate-claude} and reproduced as the Eiffel Tower Llama~\cite{eiffel-tower-llama}, the technique injects learned direction vectors into a transformer's residual stream during each forward pass.

\subsection{Steering Mechanism}

Given a steering direction $\mathbf{d} \in \mathbb{R}^{d_\text{model}}$ (a normalized SAE decoder column) and strength $\alpha$, the steered hidden state at layer $\ell$ is:

\begin{equation}
    \mathbf{h}'_\ell = \mathbf{h}_\ell + \alpha \cdot \hat{\mathbf{d}} - \text{proj}_{\hat{\mathbf{d}}}(\mathbf{h}_\ell)
    \label{eq:steering}
\end{equation}

where the projection subtraction (clamping) prevents runaway amplification when the model's existing activation already aligns with the steering direction. This is implemented as a PyTorch forward hook on the target layer---no gradient computation or weight modification required.

\subsection{Experiment Design}

We steer three layers at different depths and strengths to create non-uniform access patterns:

\begin{table}[htbp]
\centering
\caption{SAE steering configuration}
\label{tab:steering-config}
\begin{tabular}{@{}lcll@{}}
\toprule
\textbf{Layer} & \textbf{Strength ($\alpha$)} & \textbf{Position} & \textbf{Expected Effect} \\
\midrule
8  & 3.0 & Early (feature extraction) & Moderate output shift \\
16 & 5.0 & Mid (reasoning) & Strong behavior change \\
24 & 2.0 & Late (generation) & Light token distribution shift \\
\bottomrule
\end{tabular}
\end{table}

The same WFA classifier from Part~1 runs on both the unsteered baseline and the steered model, allowing direct comparison of classification outcomes.

\subsection{WFA Impact}

Steering adds a per-token forward hook overhead at steered layers. The WFA detects this as modified temporal access patterns---steered layers show different inter-access timing distributions compared to unsteered layers. In the RDMA middleware, this information drives QP allocation: steered layers with higher compute overhead benefit from dedicated low-latency Queue Pairs to avoid becoming a pipeline bottleneck.

This validates the thesis claim that SAE feature steering is a gradient-free alternative to fine-tuning (Section~11 of Update~10) and demonstrates that the WFA classifier correctly adapts its routing decisions when the inference workload changes.

%=============================================================================
\section{Next Steps}
\label{sec:next-steps}
%=============================================================================

\begin{enumerate}
    \item Run the Colab notebook on the H100 VM and capture access traces, WFA classifications, and steering comparison figures.
    \item Generate the steered vs.\ unsteered comparison figure for inclusion in Update~10.
    \item Extend the steering experiment with real SAE decoder weights (e.g., \texttt{andyrdt/saes-llama-3.1-8b-instruct}) on a Llama model.
    \item Extend the experiment to batch inference (multiple assembly samples) to produce non-uniform classification distributions.
    \item Port the validated classifier and steering engine to the two-tower cluster for integration with the RDMA middleware (Project P10).
    \item Explore fine-tuning LLM4Decompile on RISC-V assembly using the heterogeneous training pipeline (Project P16 / Section~14.2).
\end{enumerate}

\bibliographystyle{plain}
\begin{thebibliography}{1}

\bibitem{llm4decompile}
H.~Tan, Q.~Luo, J.~Li, and Y.~Zhang.
\newblock {LLM4Decompile}: Decompiling Binary Code with Large Language Models.
\newblock In \emph{Proceedings of the 2024 Conference on Empirical Methods in Natural Language Processing (EMNLP)}, 2024.
\newblock \url{https://github.com/albertan017/LLM4Decompile}

\bibitem{golden-gate-claude}
Anthropic.
\newblock Golden Gate Claude.
\newblock Blog post, May 2024.
\newblock \url{https://www.anthropic.com/news/golden-gate-claude}

\bibitem{eiffel-tower-llama}
D.~Louapre.
\newblock The Eiffel Tower Llama: Recreating Golden Gate Claude with Open-Source SAE Steering.
\newblock Hugging Face Space, 2025.
\newblock \url{https://huggingface.co/spaces/dlouapre/eiffel-tower-llama}

\end{thebibliography}

\end{document}
